\section{Proofs and Derivations}
\label{app:proofs}

\subsection{Proof of Proposition \ref{prop:nash-set}}

Equation \eqref{eq:payoff-affine} writes the payoff as
\[
W_i(x_i,x_j)
=W_i(0,x_j)+x_i\left(\Delta+\alpha A-Ax_j\right).
\]
Thus the coefficient on $x_i$ determines the global best response on $[0,1]$. It is positive, zero, or negative according as $x_j<q$, $x_j=q$, or $x_j>q$, which gives \eqref{eq:best-response}.

If $A<A^N$, then $q>1$, so every feasible rival action lies below $q$ and each player's unique best response is $1$. Hence $(1,1)$ is the unique equilibrium. If $A=A^N$, then $q=1$. A player strictly prefers $1$ against any rival action below $1$ and is indifferent over $[0,1]$ against a rival action equal to $1$. Hence a profile is an equilibrium if and only if $x_1=1$ or $x_2=1$.

Now let $A>A^N$, so $q\in(0,1)$, and consider any equilibrium. If $x_2<q$, then $x_1=1$; because $1>q$, player 2 must then choose $x_2=0$, giving $(1,0)$. If $x_2>q$, then $x_1=0$; because $0<q$, player 2 must choose $x_2=1$, giving $(0,1)$. Finally, if $x_2=q$, then $x_1<q$ would require $x_2=1$ and $x_1>q$ would require $x_2=0$, both contradictions. Thus $x_1=q$. The three profiles $(1,0)$, $(0,1)$, and $(q,q)$ all satisfy the best-response correspondence, so they are exactly the equilibria.

\subsection{Proof of Proposition \ref{prop:planner}}

Ignoring the constant $2b_D$, write
\[
\widetilde W(x_1,x_2)
=(\Delta+A)(x_1+x_2)-2Ax_1x_2.
\]
For $A<\Delta$,
\[
\widetilde W(1,1)-\widetilde W(x_1,x_2)
=(\Delta-A)(2-x_1-x_2)
+2A(1-x_1)(1-x_2)\geq0.
\]
Both terms are nonnegative on $[0,1]^2$. Equality in the first term requires $x_1+x_2=2$, which under feasibility implies $(x_1,x_2)=(1,1)$. Hence $(1,1)$ is the unique optimum.

At $A=\Delta$,
\[
\widetilde W(1,0)-\widetilde W(x_1,x_2)
=2A(1-x_1)(1-x_2)\geq0.
\]
Equality holds if and only if $x_1=1$ or $x_2=1$, giving the stated boundary optimum set.

Finally suppose $A>\Delta$ and let $s=x_1+x_2$. If $s\leq1$,
\[
\widetilde W(1,0)-\widetilde W(x_1,x_2)
=(\Delta+A)(1-s)+2Ax_1x_2\geq0.
\]
If $s\geq1$,
\[
\widetilde W(1,0)-\widetilde W(x_1,x_2)
=(A-\Delta)(s-1)+2A(1-x_1)(1-x_2)\geq0.
\]
In either branch equality requires $s=1$ and the relevant product to be zero, which yields only $(1,0)$ and $(0,1)$. These two profiles have the same welfare by symmetry, so they are the exact optima.

\subsection{Proof of Proposition \ref{prop:threshold-ordering}}

Because $0<\alpha<1$, we have $0<1-\alpha<1$. With $\Delta>0$,
\[
A^P=\Delta<\frac{\Delta}{1-\alpha}=A^N.
\]
Separately,
\[
A^N-A^P
=\frac{\Delta}{1-\alpha}-\Delta
=\frac{\alpha\Delta}{1-\alpha}>0.
\]

\subsection{Proof of Proposition \ref{prop:duplication-wedge}}

If $\Delta<A<\Delta/(1-\alpha)$, then $A>A^P$ and $A<A^N$. Proposition \ref{prop:nash-set} therefore gives the unique Nash equilibrium $(1,1)$, while Proposition \ref{prop:planner} gives the coordinated optima $(1,0)$ and $(0,1)$. Moreover,
\[
W(1,0)-W(1,1)=A-\Delta>0,
\]
so coordinated reallocation is strictly welfare improving.

\subsection{The symmetric high-$A$ equilibrium}

For $A>A^N$, $q\in(0,1)$. Substituting $\Delta=A(q-\alpha)$ into \eqref{eq:aggregate-welfare} gives
\[
W(1,0)-W(q,q)
=A\left[(1-\alpha)(1-q)+\alpha q\right].
\]
Here $1-\alpha>0$, $1-q>0$, $\alpha>0$, and $q>0$, so both terms inside the brackets are strictly positive. This proves \eqref{eq:highA-gap}.

\subsection{Proof of Proposition \ref{prop:switching-lemma}}

Relative to duplication, the coordinated net gain from reallocation is
\[
G-\Delta,
\]
whereas the reallocating jurisdiction's net gain is
\[
\lambda G-\Delta.
\]
The planner therefore switches exactly when $G>\Delta$, while the local government switches exactly when $G>\Delta/\lambda$. Thus
\[
G^P=\Delta,
\qquad
G^N=\frac{\Delta}{\lambda}.
\]
Since $0<\lambda<1$ and $\Delta>0$, $G^P<G^N$.

\subsection{Integration thresholds and endogenous capture}

Under the frozen integration assumptions, $G$ is continuous and strictly increasing, $\lambda\in(0,1)$ is constant, and both $\Delta$ and $\Delta/\lambda$ lie in the range of $G$. Hence there are unique crossings
\[
G(\tau^P)=\Delta,
\qquad
G(\tau^N)=\frac{\Delta}{\lambda}.
\]
For their ordering, only strict monotonicity is needed: because $\Delta<\Delta/\lambda$, we obtain $\tau^P<\tau^N$. Thus continuity and the range conditions supply existence, while strict monotonicity supplies uniqueness and ordering.

For the endogenous-capture example, let $G(\tau)=1+\tau$, $\lambda(\tau)=(1+\tau)^{-2}$, and $\Delta=1$. Then
\[
\lambda(\tau)G(\tau)=\frac{1}{1+\tau}.
\]
Thus $G(\tau)>1$ for every $\tau>0$, while $\lambda(\tau)G(\tau)\leq1$ for every $\tau\geq0$. The planner therefore switches for positive $\tau$, but the local government never has a strict switching incentive. Incomplete capture alone does not guarantee a finite decentralized switching threshold.

\subsection{Restricted constant-returns production result}

Let $F$ satisfy the conditions in Section \ref{sec:robustness}. Degree-one homogeneity and differentiability imply Euler's identity
\[
F(1,1)=F_u(1,1)+F_d(1,1).
\]
Since $F_u(1,1)>0$ and $F_d(1,1)>0$,
\[
F(1,1)>F_d(1,1)>0.
\]
The coordinated and local thresholds are respectively
\[
A_F^P=\frac{\Delta}{F(1,1)},
\qquad
A_F^N=\frac{\Delta}{F_d(1,1)},
\]
so $A_F^P<A_F^N$.

\subsection{CES specialization}

For
\[
F(u,d)=\left[\omega u^\rho+(1-\omega)d^\rho\right]^{1/\rho},
\]
the downstream marginal product is
\[
F_d(u,d)
=(1-\omega)d^{\rho-1}
\left[\omega u^\rho+(1-\omega)d^\rho\right]^{1/\rho-1}.
\]
At the unit match,
\[
F(1,1)=1,
\qquad
F_d(1,1)=1-\omega.
\]
Hence $A_F^P=\Delta$ and $A_F^N=\Delta/(1-\omega)$ on the admissible concave CES parameter domain.

\subsection{Alternative matching identity and local best response}

Let $s=x_1+x_2$. From the two minimum terms in \eqref{eq:capacity-matching},
\[
M(x_1,x_2)
=\begin{cases}
s, & s\leq1,\\
2-s, & s\geq1,
\end{cases}
=1-|s-1|.
\]
Fix rival allocation $x_j=y$. Region $i$'s nonconstant payoff is
\[
\Delta x_i+(1-\lambda)A\min\{x_i,1-y\}
+\lambda A\min\{1-x_i,y\}.
\]
Its slope is $\Delta+(1-\lambda)A>0$ for $x_i\leq1-y$ and $\Delta-\lambda A$ for $x_i\geq1-y$. Therefore, if $A<\Delta/\lambda$, the payoff is strictly increasing over $[0,1]$, so $x_i=1$ is the unique best response to every $y$. Hence $(1,1)$ is the unique decentralized equilibrium.

For the planner,
\[
\Delta s+A M
=\begin{cases}
(\Delta+A)s, & s\leq1,\\
2A+(\Delta-A)s, & s\geq1.
\end{cases}
\]
When $A>\Delta$, the first branch is strictly increasing and the second strictly decreasing, so the exact coordinated optimum set is $s=1$. Combining the two conditions gives the robustness wedge
\[
\Delta<A<\frac{\Delta}{\lambda}.
\]
This establishes reallocation away from duplication under the stated alternative matching rule, without implying complete regional specialization.
